% !TEX program = xelatex
% ===========================================================================
%  因式分解专题讲义—— 中文数学习题课讲义（Loom 模板 · 课前自学版）
%  必须用 XeLaTeX 编译两遍；QA 用 scripts/qa_check.py。
%  第一行的魔法注释让 VS Code (LaTeX Workshop) 走 XeLaTeX——别删（默认是
%  pdfLaTeX，会报 xeCJK requires XeTeX）。
% ===========================================================================
\documentclass{loom}

% 挑战题 ★ 用 Zapf Dingbats 绘制，不依赖系统字体（Optima 等）是否含 U+2605
\usepackage{pifont}
\usetikzlibrary{decorations.pathreplacing}

% 页眉线（可选）；老师偏好极简时可整行注释掉
%\runningthread{因式分解专题讲义 · 课前自学讲义}

% ---------------------------------------------------------------------------
%  讲义专用组件：题面（茜红结）、解析（靛蓝结）、三段式解答标签
%  —— 这些宏定义住在 main.tex（不在 cls），以便随讲义微调。
% ---------------------------------------------------------------------------
% 题目框：breakable=false 保证一题不跨页。默认不留答题空白区（老师偏好）。
% —— 若要做「打印作答版」，取消下面 \answerspace 定义与 problem 框 after upper
%    一行的注释即可，题末会追加一块纯白书写面板（默认 5 行，可 \answerspace[8]）。
% \newcommand{\answerspace}[1][5]{\par\medskip\noindent
%   \colorbox{white}{\parbox[t][#1\baselineskip][t]{\dimexpr\linewidth-2\fboxsep\relax}{\strut}}}
\newtcolorbox{problem}[1]{knotmadder, breakable=false, before skip=6pt plus 2pt, after skip=6pt plus 2pt,
  before upper={{\headingfont\bfseries\color{madder}#1}\par\smallskip}%
  % , after upper={\answerspace}   % 打印作答版：取消本行注释
}
\newtcolorbox{solution}[2][]{knotindigo, lines before break=4,
  before skip=6pt plus 2pt, after skip=6pt plus 2pt,
  before upper={{\headingfont\bfseries\color{indigo}#2}\par\smallskip},#1}

\newcommand{\think}{\par\smallskip\noindent
  {\headingfont\bfseries\color{indigo}〔怎么想〕}\enspace}
\newcommand{\solve}{\par\smallskip\noindent
  {\headingfont\bfseries\color{inkiron}〔解〕}\enspace}
\newcommand{\solvehead}[1]{\par\smallskip\noindent
  {\headingfont\bfseries\color{inkiron}〔#1〕}\enspace}
\newcommand{\insight}[1]{\par\nopagebreak[4]\smallskip\nopagebreak[4]\noindent
  {\headingfont\bfseries\color{weld!60!inkiron}〔方法点睛〕}\enspace
  {\strandfont\small #1}\par}
\newcommand{\hint}[1]{\par\smallskip\noindent
  {\strandfont\small{\raisebox{-0.8pt}{\loomtile{weld}}\hspace{4pt}}%
  {\bfseries\color{weld!60!inkiron}入手提示}\enspace{\color{inkiron!85!white}#1}}\par}
% 页边旁批：放"想一想/追问"、∑/∏ 记号说明，不打断正文
\newcommand{\sidenote}[1]{\marginnote{\raggedright\strandfont\footnotesize
  \color{selvage}#1}}
\newcommand{\starred}{{\color{weld!70!inkiron}\ding{72}}}

\setlist[enumerate]{leftmargin=2.3em, itemsep=2pt, topsep=3pt}
\setlist[itemize]{leftmargin=1.6em, itemsep=2pt, topsep=3pt}


% 例题随文讲解；练习与选学挑战的解析集中放在书末。
\newif\ifstudentnotes
\studentnotesfalse
\let\originalhint\hint
\renewcommand{\hint}[1]{\ifstudentnotes\else\originalhint{#1}\fi}
\newcommand{\methodcard}[2]{\ifstudentnotes\else
\begin{tcolorbox}[knotindigo,breakable=false,before skip=6pt plus 2pt,after skip=6pt plus 2pt]
\textbf{方法小结}\par #1\par\smallskip
\textbf{易错提醒}\quad #2\end{tcolorbox}\fi}
\newcommand{\topicline}[1]{\ifstudentnotes\else\noindent\whisper{#1}\par\medskip\fi}
\setcounter{secnumdepth}{2}
\setcounter{tocdepth}{2}
\renewcommand{\contentsname}{目录}
\newcommand{\chapterordinal}{\ifcase\value{section}\or 一\or 二\else\arabic{section}\fi}
\titleformat{\section}[block]{\sectionheadfont}{第\chapterordinal 章}{0.7em}{}
  [\vspace{1pt}{\color{madder}\rule{\linewidth}{1.1pt}}]
\titleformat{\subsection}[block]{\headingfont\bfseries\color{indigo}\fontsize{12}{15}\selectfont}
  {\thesubsection}{0.65em}{}
\titlespacing*{\subsection}{0pt}{16pt plus 3pt}{7pt}
\hypersetup{pdftitle={因式分解专题讲义},pdfauthor={数学专题讲义}}

% 封面：四块面积相加，等于矩形两边的乘积。
\definecolor{tilegreen}{HTML}{2E806D}
\newcommand{\factorcoverfigure}{%
\begin{tikzpicture}[x=0.9cm,y=0.9cm,line width=0.6pt,
  every node/.style={text=inkiron,font=\normalsize}]
  \fill[indigo!23!white] (0,0) rectangle (2.55,2.55);
  \fill[weld!40!white] (2.55,0) rectangle (3.65,2.55);
  \fill[madder!24!white] (0,2.55) rectangle (2.55,3.45);
  \fill[tilegreen!32!white] (2.55,2.55) rectangle (3.65,3.45);
  \draw[inkiron!65] (0,0) rectangle (3.65,3.45);
  \draw[inkiron!45] (2.55,0)--(2.55,3.45) (0,2.55)--(3.65,2.55);
  \node[font=\large,text=indigo] at (1.275,1.275) {$x^2$};
  \node[font=\large,text=weld!65!inkiron] at (3.10,1.275) {$px$};
  \node[font=\large,text=madder] at (1.275,3.0) {$qx$};
  \node[font=\large,text=tilegreen!75!black] at (3.10,3.0) {$pq$};
  \node[font=\large] at (1.825,-0.65)
    {$\textcolor{indigo}{x^2}+(p+q)x+\textcolor{tilegreen!75!black}{pq}=(x+p)(x+q)$};
\end{tikzpicture}}
% 只替换封面主图；讲义中的结纹图标仍使用原来的 \loomtile。
\renewcommand{\weaveemblem}[1]{\factorcoverfigure}

\newcommand{\factorfigure}{%
\begin{center}
\begin{tikzpicture}[scale=0.78,line width=0.6pt]
\filldraw[fill=indigo!12] (0,0) rectangle (1.05,1.05);
\filldraw[fill=indigo!12] (1.05,0) rectangle (2.55,1.5);
\node[below] at (0.525,0) {$a$};
\node[left] at (0,0.525) {$a$};
\node[below] at (1.8,0) {$b$};
\node[right] at (2.55,0.75) {$b$};
\node at (1.2,-0.75) {图①};
\begin{scope}[xshift=4.4cm]
\coordinate (A) at (0,2);
\coordinate (B) at (0,0);
\coordinate (P) at (2,0);
\coordinate (C) at (3,0);
\coordinate (D) at (3,1);
\fill[indigo!18] (A)--(P)--(D)--cycle;
\draw (A)--(B)--(C)--(D)--cycle;
\draw (A)--(P)--(D);
\draw[dashed,selvage] (0,2)--(2,2)--(2,0);
\draw[dashed,selvage] (2,0)--(2,1)--(3,1);
\draw (0,0.15)--(0.15,0.15)--(0.15,0);
\draw (2.85,0)--(2.85,0.15)--(3,0.15);
\node[above left] at (A) {$A$};
\node[below left] at (B) {$B$};
\node[below] at (P) {$P$};
\node[below right] at (C) {$C$};
\node[above right] at (D) {$D$};
\node[left] at (0,1) {$m$};
\node[below] at (1,-0.35) {$m$};
\node[below] at (2.5,-0.35) {$n$};
\node[right] at (3,0.5) {$n$};
\node at (1.5,-0.75) {图②};
\end{scope}
\end{tikzpicture}
\end{center}}

\begin{document}
\loomcover{因式分解专题讲义}{}{}{2026 年 10 月}
\clearpage
{\small\tableofcontents}
\clearpage
\section{因式分解的方法拓展}

\noindent\whisper{未特别说明时，本章在有理系数范围内分解因式。}\par

\Needspace{18\baselineskip}

\subsection{分组与拆项：让共同因式出现}

\topicline{前两题从已学的提公因式法出发，后两题分别接上公式法与有目标的拆项。}

\Needspace{12\baselineskip}

\begin{problem}{例题 1 （分组后再提一次）}

分解因式：$x^2-xy+4x-4y$。

\end{problem}

\Needspace{6\baselineskip}

\begin{solution}{例题 1 · 讲解}

\think 前两项都含 $x$，后两项都含 $4$；两组里 $x$ 与 $y$ 的系数比例相同，因此分组后有机会出现同一个括号。

\solve \[x^2-xy+4x-4y=x(x-y)+4(x-y)=(x-y)(x+4).\]

\insight{分组的依据是各组能产生共同因式；分组本身只是中间步骤。}

\end{solution}

\Needspace{12\baselineskip}

\begin{problem}{例题 2 （一组先配成平方）}

分解因式：$x^2-4y^2-z^2+4yz$。

\end{problem}

\Needspace{6\baselineskip}

\begin{solution}{例题 2 · 讲解}

\think $4y^2$ 与 $z^2$ 是平方，$4yz$ 是 $2\cdot2y\cdot z$。把这三项连同符号整理，可使原式成为两个平方的差。

\solve \begin{align*}
x^2-4y^2-z^2+4yz&=x^2-(4y^2-4yz+z^2)\\
&=x^2-(2y-z)^2\\
&=(x-2y+z)(x+2y-z).
\end{align*}

\insight{括号前有负号时，先检查括号内每一项的符号，再识别完全平方。}

\end{solution}

\Needspace{12\baselineskip}

\begin{problem}{例题 3 （拆项的目标是什么）}

分解因式：$x^3+6x^2+11x+6$。尝试给出两种拆项方式。

\end{problem}

\Needspace{6\baselineskip}

\begin{solution}{例题 3 · 讲解}

\think 如果想提出 $x+1$，三次项应与 $x^2$ 配成一组，常数 $6$ 应与 $6x$ 配成一组；余下 $5x^2+5x$ 正好也含 $x+1$。第二种拆法同样围绕目标因式安排系数。

\solve 第一种拆法：
\begin{align*}
x^3+6x^2+11x+6&=(x^3+x^2)+(5x^2+5x)+(6x+6)\\
&=(x+1)(x^2+5x+6)\\
&=(x+1)(x+2)(x+3).
\end{align*}
第二种拆法：
\begin{align*}
x^3+6x^2+11x+6&=(x^3+2x^2)+(4x^2+8x)+(3x+6)\\
&=(x+2)(x^2+4x+3)\\
&=(x+1)(x+2)(x+3).
\end{align*}

\insight{先确定希望出现的括号，再拆项；拆出的各项相加必须恢复原式。}

\end{solution}

\methodcard{先观察各项共有的字母或系数比例，再尝试分组。每组可以先提公因式，也可以先用公式；关键是各组处理后，整个式子还能分解。拆项时，围绕希望出现的因式安排系数。}{提出负号时，括号内各项都要变号；拆项前后逐项核对，防止少项或多项。}

\begin{problem}{练习 1 （分组以后仍要检查）}

分解因式：$a^3+a^2-a-1$。

\end{problem}

\Needspace{18\baselineskip}

\subsection{配方与添项：把结构改成平方差}

\topicline{由现成的完全平方，过渡到给两组配方，再学习添减同一项。}

\Needspace{12\baselineskip}

\begin{problem}{例题 4 （把一次项藏进平方）}

分解因式：$x^2-y^2-4x+4$。

\end{problem}

\Needspace{6\baselineskip}

\begin{solution}{例题 4 · 讲解}

\think $x^2-4x+4$ 恰好是完全平方，剩下的是 $-y^2$，因此先配方再用平方差。

\solve \[x^2-y^2-4x+4=(x-2)^2-y^2=(x-2-y)(x-2+y).\]

\insight{不要只盯着原式的项数；重新组合后可能直接显出公式。}

\end{solution}

\Needspace{12\baselineskip}

\begin{problem}{例题 5 （从四次式找平方差）}

分解因式：$x^4-7x^2+1$。

\end{problem}

\Needspace{6\baselineskip}

\begin{solution}{例题 5 · 讲解}

\think 首尾为 $x^4$ 与 $1$，可以先考虑 $(x^2+1)^2$。它的中间项是 $2x^2$，比原式多 $9x^2$，多出的部分也是平方。

\solve \begin{align*}
x^4-7x^2+1&=x^4+2x^2+1-9x^2\\
&=(x^2+1)^2-(3x)^2\\
&=(x^2-3x+1)(x^2+3x+1).
\end{align*}

\insight{添项与减项要成对出现；先计算补出的中间项与原中间项相差多少。}

\end{solution}

\methodcard{先找平方项，再用 $2AB$ 判断缺少的中间项。已经存在完全平方就先重组；需要添项时，同时减去同一项，使原式恒等不变。}{$A^2+B^2$ 不能直接套平方差；补出的中间项必须按 $2AB$ 计算。}

\begin{problem}{练习 2 （常数拆给两个平方）}

分解因式：$4x^2-4x-y^2+4y-3$。

\end{problem}

\begin{problem}{练习 3 （平方和也可能继续分解）}

在有理系数范围内分解因式：$a^4+4$。

\hint{若先写出 $(a^2+2)^2$，需要减去什么才能恢复原式？}

\end{problem}

\Needspace{18\baselineskip}

\subsection{整体与换元：减少结构的复杂程度}

\topicline{从直接重复的块开始，逐步学习先配对制造整体，最后检查代回后的各因式。}

\Needspace{12\baselineskip}

\begin{problem}{例题 6 （反复出现的块）}

分解因式：$(x^2+5x+2)(x^2+5x+3)-12$。

\end{problem}

\Needspace{6\baselineskip}

\begin{solution}{例题 6 · 讲解}

\think 两个括号都含 $x^2+5x$，将它记为 $u$，只需处理关于 $u$ 的二次三项式。代回后还要检查各因式。

\solve 令 $u=x^2+5x$，则
\[(u+2)(u+3)-12=u^2+5u-6=(u+6)(u-1).\]
代回得
\begin{align*}
\text{原式}&=(x^2+5x+6)(x^2+5x-1)\\
&=(x+2)(x+3)(x^2+5x-1).
\end{align*}
$x^2+5x-1$ 不能分解为有理系数的一次因式，本题保留该二次因式。

\insight{整体出现两次就可能值得换元；是否换元取决于结构能否简化。}

\end{solution}

\Needspace{12\baselineskip}

\begin{problem}{例题 7 （换元之后还有一层）}

分解因式：$(x^4+x^2-4)(x^4+x^2+3)+10$。

\end{problem}

\Needspace{6\baselineskip}

\begin{solution}{例题 7 · 讲解}

\think 换元能先降低结构的复杂程度。代回后，一个因式有平方差，另一个因式可以添减 $x^2$，因此不能停在两个高次因式的乘积。

\solve 令 $u=x^4+x^2$，则
\[(u-4)(u+3)+10=u^2-u-2=(u-2)(u+1).\]
其中
\begin{align*}
x^4+x^2-2&=(x^2-1)(x^2+2)=(x-1)(x+1)(x^2+2),\\
x^4+x^2+1&=(x^2+1)^2-x^2=(x^2-x+1)(x^2+x+1).
\end{align*}
故原式为
\[(x-1)(x+1)(x^2+2)(x^2-x+1)(x^2+x+1).\]

\insight{高次因式要继续看结构；“看起来像平方和”不等于不能分解。}

\end{solution}

\methodcard{圈出相同部分；必要时先重组使它出现；用新字母简化；用已学方法分解；换回原字母并继续检查。整体出现的次数没有固定门槛。}{换元后必须代回；得到完全平方或高次因式后仍须检查分解是否彻底。}

\begin{problem}{练习 4 （先配对，再找整体）}

分解因式：$(x+1)(x+2)(x+3)(x+6)+x^2$。

\hint{比较不同配对方式相乘后的一次项和常数项。}

\end{problem}

\begin{problem}{练习 5 （两个整体合成一个）}

分解因式：$(x+y)(x+y+2xy)+(xy+1)(xy-1)$。

\end{problem}

\begin{problem}{练习 6 （完全平方以后继续分解）}

分解因式：$(x+y-2xy)(x+y-2)+(xy-1)^2$。

\end{problem}

\Needspace{18\baselineskip}

\subsection{主元整理：选一个字母建立秩序}

\topicline{先整理二元二次式，再换主元暴露公因式，最后处理循环结构。}

\Needspace{12\baselineskip}

\begin{problem}{例题 8 （把另一个字母当系数）}

分解因式：$x^2+xy-6y^2+x+13y-6$。

\end{problem}

\Needspace{6\baselineskip}

\begin{solution}{例题 8 · 讲解}

\think 把 $y$ 暂时视作常数，原式是关于 $x$ 的二次三项式。常数部分 $-6y^2+13y-6$ 可以先分解，再检查两个括号的一次项和。

\solve \begin{align*}
\text{原式}&=x^2+(y+1)x-(3y-2)(2y-3)\\
&=[x+(3y-2)][x-(2y-3)]\\
&=(x+3y-2)(x-2y+3).
\end{align*}
两个括号中常数部分相加为 $(3y-2)+(-2y+3)=y+1$，与 $x$ 的系数一致。

\insight{选一个字母整理成熟悉的二次式，其他字母就成为系数。}

\end{solution}

\Needspace{12\baselineskip}

\begin{problem}{例题 9 （循环式里寻找相同括号）}

分解因式：$a^2(b-c)+b^2(c-a)+c^2(a-b)$。

\end{problem}

\Needspace{6\baselineskip}

\begin{solution}{例题 9 · 讲解}

\think 把 $a$ 当主元后，二次项系数为 $b-c$，一次项系数为 $c^2-b^2$，常数项为 $b^2c-bc^2$。后三者都可以提取 $b-c$。

\solve \begin{align*}
\text{原式}&=(b-c)a^2+(c^2-b^2)a+bc(b-c)\\
&=(b-c)[a^2-(b+c)a+bc]\\
&=(b-c)(a-b)(a-c).
\end{align*}
进一步，若原式等于零，则 $(b-c)(a-b)(a-c)=0$，至少一个差为零，因此 $a,b,c$ 中至少有两个相等。

\insight{循环排列的式子可以先打破排列，按一个字母整理出共同结构。}

\end{solution}

\methodcard{选一个字母，按它的降幂收集各项，把其余字母暂作系数。可以先试次数较低、系数较易分解的字母；若整理没有帮助，再换一个角度。}{收集同次项时保留所有系数与符号；主元的选择是观察线索，不是固定结论。}

\begin{problem}{练习 7 （换一个主元看公因式）}

分解因式：$x^3+(2a+1)x^2+(a^2+2a-1)x+(a^2-1)$。

\hint{比较 $x$ 与 $a$ 的最高次数；尝试按 $a$ 的降幂整理。}

\end{problem}

\Needspace{18\baselineskip}

\subsection{比较系数：由目标因式倒推}

\topicline{从确定常数项，到多元参数，再到已知二次因式；试代值仅作为拆项线索。}

\Needspace{12\baselineskip}

\begin{problem}{例题 10 （用常数项确定第三个因式）}

已知 $x^3+ax^2+bx+8$ 有两个因式 $x+1$ 和 $x+2$，求 $a+b$。

\end{problem}

\Needspace{6\baselineskip}

\begin{solution}{例题 10 · 讲解}

\think 原式首项系数为 $1$，已经知道两个不同的一次因式，第三个因式应写成 $x+c$。先比较常数项，能直接确定 $c$。

\solve 设
\[x^3+ax^2+bx+8=(x+1)(x+2)(x+c).\]
比较常数项，$2c=8$，得 $c=4$。于是
\[(x+1)(x+2)(x+4)=x^3+7x^2+14x+8,\]
所以 $a=7$，$b=14$，$a+b=21$。

\insight{比较系数时，先选最容易确定的常数项或最高次项。}

\end{solution}

\Needspace{12\baselineskip}

\begin{problem}{例题 11 （由一次项倒推参数）}

$k$ 为何值时，多项式 $x^2-2xy+ky^2+3x-5y+2$ 能分解为两个实系数的一次因式的积？

\end{problem}

\Needspace{6\baselineskip}

\begin{solution}{例题 11 · 讲解}

\think 令 $y=0$ 得 $(x+1)(x+2)$。因此两个一次因式经调整常数倍后，可分别设成 $x+my+1$ 与 $x+ny+2$。比较 $xy$ 和 $y$ 的系数，先确定 $m,n$。

\solve 设原式为 $(x+my+1)(x+ny+2)$，展开并比较系数得
\[m+n=-2,\qquad 2m+n=-5,\qquad mn=k.\]
前两式相减得 $m=-3$，继而 $n=1$，所以 $k=-3$。
此时原式确为 $(x-3y+1)(x+y+2)$，满足要求。

\insight{先利用已知的简单部分确定因式形状，再用系数条件求参数。}

\end{solution}

\methodcard{根据次数与首项系数设出待求因式；展开；比较同次项系数；从最简单的等式开始求参数；最后检查所得积确为原式。}{题中“有公因式”应指非恒定公因式；本节不用长除法、二次公式或判别式。}

\begin{problem}{练习 8 （用拆项落实候选因式）}

分解因式：$x^3-2x^2-13x-10$。

\hint{代入 $x=-1$ 后式子为零；尝试通过拆项让每组都出现 $x+1$。}

\end{problem}

\begin{problem}{练习 9 （已知二次因式，比较四次式系数）}

已知 $x^2+x-6$ 是多项式 $2x^4+x^3-ax^2+bx+a+b-1$ 的因式，求 $a,b$。

\hint{把另一个因式设为二次式；先比较最高次项。}

\end{problem}

\Needspace{12\baselineskip}

\subsection{巩固与选学}

\topicline{前两题检查是否会独立识别结构；最后一题选学，不作为基本达标要求。}

\begin{problem}{练习 10 （巩固：识别一个平方）}

分解因式：$x^4+2x^3+3x^2+2x+1$。

\end{problem}

\methodcard{先说清观察依据，再写变形。遇到难题时，先寻找式子各部分的联系，随后决定使用分组、配方或整体思想。}{本章默认在有理系数范围内分解；不要求为拆二次式引入根式系数。}

\begin{problem}{练习 11 （巩固：分组方式需要改变）}

在实数范围内分解因式：$x^2y^2-x^2y-xy^2+x+y-1$。

\end{problem}

\begin{problem}{挑战 1 \starred （选学：三个整体之间的联系）}

分解因式：$(2x-3y)^3+(3x-2y)^3-125(x-y)^3$。本题可使用立方和公式。

\hint{比较前两个立方底数的和与第三个立方底数。}

\end{problem}


\Needspace{24\baselineskip}
\section{因式分解的应用}

\Needspace{18\baselineskip}

\subsection{条件求值}

\topicline{题目只求某个整体时，先看已知条件能否直接确定它。}

\Needspace{12\baselineskip}

\begin{problem}{例题 15 （恒等式中的参数）}

设 $y=ax$，若代数式 $(x+y)(x-2y)+3y(x+y)$ 可化简为 $x^2$，求 $a$。这里“化简为”表示关于 $x$ 的恒等式。

\end{problem}

\Needspace{6\baselineskip}

\begin{solution}{例题 15 · 讲解}

\think 两项都有 $x+y$，先提出它，另一个因式也恰好是 $x+y$。代入 $y=ax$ 后，比较 $x^2$ 的系数即可，不必分别求 $x,y$。

\solve \begin{align*}
(x+y)(x-2y)+3y(x+y)&=(x+y)(x+y)\\
&=(1+a)^2x^2.
\end{align*}
由恒等式，比较系数得 $(1+a)^2=1$，即 $1+a=\pm1$，所以 $a=0$ 或 $a=-2$。
代入 $a=0$ 或 $a=-2$，原式的系数 $(1+a)^2$ 都为 $1$。

\insight{恒等式可比较系数；解平方等式时要保留两种符号。}

\end{solution}

\Needspace{12\baselineskip}

\begin{problem}{例题 16 （由两个条件求和）}

实数 $x,y$ 满足 $x^2+xy+y=14$、$y^2+xy+x=28$，求 $x+y$。

\end{problem}

\Needspace{6\baselineskip}

\begin{solution}{例题 16 · 讲解}

\think 题目只问 $x+y$，两式相加就能凑出它的平方。求出的两个和值都要保留，再用两式之差分别求出一组 $x,y$，确认原条件确实允许这些和值。

\solve 令 $t=x+y$，两式相加得
\[t^2+t=42,\qquad (t-6)(t+7)=0.\]
故 $t=6$ 或 $t=-7$。两式相减得 $(x-y)(t-1)=-14$。
当 $t=6$ 时，$x-y=-\frac{14}{5}$，得 $(x,y)=(\frac85,\frac{22}5)$；
当 $t=-7$ 时，$x-y=\frac74$，得 $(x,y)=(-\frac{21}8,-\frac{35}8)$。
将两组数分别代回原条件均成立，所以 $x+y=6$ 或 $-7$。

\insight{求出整体的候选值后，要确认它们都能满足原条件。}

\end{solution}

\methodcard{先看条件与目标式中有没有相同的整体。条件相加、相减或移项后，有时能直接求出这个整体，无须分别求每个字母。}{换元后保留原有范围；题目未说字母为正时，要保留负值的可能。}

\begin{problem}{练习 17}

实数 $a,b$ 满足 $a^4+b^4=a^2-2a^2b^2+b^2+6$，求 $a^2+b^2$。

\end{problem}

\begin{problem}{练习 18}

两个不等实数 $m,n$ 满足 $m^2-2m=a$、$n^2-2n=a$、$m^2+n^2=5$，求 $a$。

\end{problem}

\Needspace{18\baselineskip}

\subsection{简便计算}

\topicline{数字较大或因数较多时，先看各项的关系，再决定怎样分解。}

\Needspace{12\baselineskip}

\begin{problem}{例题 17 （大数计算）}

计算 $\displaystyle\frac{2004^3-2\times2004^2-2002}{2004^3+2004^2-2005}$。

\end{problem}

\Needspace{6\baselineskip}

\begin{solution}{例题 17 · 讲解}

\think 这些数都接近 $2004$。令 $t=2004$，则 $2002=t-2$、$2005=t+1$，分子、分母都能写成四项式，再用分组法分解。

\solve 令 $t=2004$，则
\begin{align*}
\text{分子}&=t^3-2t^2-t+2=(t-2)(t^2-1),\\
\text{分母}&=t^3+t^2-t-1=(t+1)(t^2-1).
\end{align*}
因 $t=2004$，上述分母和被约去的因式都不为零，故
\[\text{原式}=\frac{t-2}{t+1}=\frac{2002}{2005}.\]

\insight{大数之间有简单的和差关系时，先用字母表示，往往更容易计算。}

\end{solution}

\Needspace{12\baselineskip}

\begin{problem}{例题 18 （连乘计算）}

计算
\[
\frac{(2\times5+2)(4\times7+2)\cdots(1994\times1997+2)}
{(1\times4+2)(3\times6+2)\cdots(1993\times1996+2)}.
\]

\end{problem}

\Needspace{6\baselineskip}

\begin{solution}{例题 18 · 讲解}

\think 先看第一项：$2\times5+2=3\times4$，$1\times4+2=2\times3$，相除得到 $2$。第二组会得到 $3/2$ 吗？用第 $k$ 组的式子验证，就能看清后面的约分规律。

\solve 分子、分母各有 $997$ 项。对 $k=1,2,\ldots,997$，有
\begin{align*}
2k(2k+3)+2&=2(2k+1)(k+1),\\
(2k-1)(2k+2)+2&=2k(2k+1).
\end{align*}
所以
\[\text{原式}=\frac21\times\frac32\times\frac43\times\cdots\times\frac{998}{997}=998.\]

\insight{先分解一组对应项，再写一般项，最后核对首项、末项和项数。}

\end{solution}

\methodcard{先用字母写出数字间的关系。连乘题可从一组对应项入手，分解后再看哪些因式重复出现。}{只有分子、分母中的相同非零因式才能约去，和式中的项不能直接约。}

\begin{problem}{练习 19}

计算
\[
\frac{(7^4+64)(15^4+64)(23^4+64)(31^4+64)(39^4+64)}
{(3^4+64)(11^4+64)(19^4+64)(27^4+64)(35^4+64)}.
\]

\end{problem}

\Needspace{18\baselineskip}

\subsection{平方与非负性}

\topicline{配成平方后，除了可以分解，还可以利用平方非负判断取值。}

\Needspace{12\baselineskip}

\begin{problem}{例题 19 （平方和为零）}

已知实数 $m,n$，当 $x=m$ 时，多项式 $x^2+2x+n^2$ 的值为 $-1$；求 $x=-m$ 时该多项式的值。

\end{problem}

\Needspace{6\baselineskip}

\begin{solution}{例题 19 · 讲解}

\think 把 $x=m$ 代入，得 $m^2+2m+n^2=-1$。移项后能配成两个平方的和；和为零时，两个平方都只能为零。

\solve 由已知条件得
\[(m+1)^2+n^2=0.\]
两项均非负，故 $m+1=0$、$n=0$，即 $m=-1$。
当 $x=-m=1$ 时，原多项式的值为 $1^2+2\times1+0=3$。

\insight{实数的平方非负；几个平方之和为零时，每个平方都等于零。}

\end{solution}

\Needspace{12\baselineskip}

\begin{problem}{例题 20 （和一定，求积的最大值）}

观察 $1\times49,2\times48,\ldots,25\times25,\ldots,49\times1$。
\begin{enumerate}[label=(\arabic*)]
\item 写出两数乘积的最大值，并写出两个因数 $a,b$ 的关系。
\item 类比 $1\times59,2\times58,\ldots,59\times1$，求乘积最大值并证明。
\end{enumerate}

\end{problem}

\Needspace{6\baselineskip}

\begin{solution}{例题 20 · 讲解}

\think 第一组两因数的和总是 $50$，乘积在两数接近时较大。怎样证明？由 $4ab=(a+b)^2-(a-b)^2$，和固定后，差的平方越小，乘积越大。也可以由恒等式 $(a-b)^2=(a+b)^2-4ab\ge0$，得到 $(a+b)^2\ge4ab$。

\solve 第一组有 $a+b=50$，所以
\[ab=\frac{2500-(a-b)^2}{4}\le625.\]
当 $a=b=25$ 时等号成立，故最大值为 $625$。
第二组记两因数为 $m,n$，则 $m+n=60$，于是
\[mn=\frac{3600-(m-n)^2}{4}\le900.\]
当 $m=n=30$ 时等号成立，故最大值为 $900$。

再想一步：若两个正整数的和为 $51$，它们不可能相等，差的绝对值最小是多少？

\insight{证明最大值时，既要给出上界，也要说明等号能够取到。}

\end{solution}

\methodcard{可从交叉项决定怎样配方。配成平方后，要接着说明它对取值有什么限制：平方和为零、平方非负，或某个平方有上界。}{和为零不能直接推出每项为零。求最值时还要检查等号条件。}

\begin{problem}{练习 20}

实数 $x,y$ 满足 $x^2+12xy+52y^2-8y+1=0$，求 $x^2-y^2$。

\end{problem}

\begin{problem}{练习 21}

已知实数 $m\ne n$，$a=2m^2-mn$、$b=mn-2n^2$、$c=m^2-n^2$。用“$<$”表示 $a,b,c$ 的大小关系。

\end{problem}

\begin{problem}{练习 22}

有多少个整数解 $(x,y)$ 满足 $x^2+y^2\le2x+2y$？

\end{problem}

\Needspace{18\baselineskip}

\subsection{整数方程}

\topicline{把方程配成两个整数的乘积，常能将求解变为列举有限个因数对。}

\Needspace{12\baselineskip}

\begin{problem}{例题 21 （整数解的列举）}

求方程 $xy-2x-2y+7=0$ 的全部整数解，其中 $x\le y$。

\end{problem}

\Needspace{6\baselineskip}

\begin{solution}{例题 21 · 讲解}

\think 先把 $xy-2x-2y$ 与 $(x-2)(y-2)$ 比较，后者只多一个常数 $4$。补上这项，方程就成了两个整数的乘积等于 $-3$，可列出因数对。

\solve 原方程等价于
\[xy-2x-2y+4=-3,\qquad (x-2)(y-2)=-3.\]
整数因数对在 $x-2\le y-2$ 的限制下只有 $(-3,1)$、$(-1,3)$。
因此 $(x,y)=(-1,3)$ 或 $(1,5)$。

\insight{配成整数乘积后，先列因数对，再用大小条件筛选。}

\end{solution}

\Needspace{6\baselineskip}

\methodcard{含有 $xy$ 及 $x,y$ 一次项的方程，可尝试添一个常数配成乘积。例如 $xy+x+y=M$ 可写成 $(x+1)(y+1)=M+1$（回看练习 7、8 后的总结）。若 $x,y$ 为正整数，新因数都至少为 $2$；这些范围要一并用来筛选。}{整数可以为负。}

\begin{problem}{练习 23}

整数 $a,b$ 满足 $6ab=9a-10b+303$，求 $a+b$。

\end{problem}

\begin{problem}{练习 24}

从 $1,2,\ldots,37$ 中选出两个不同的数，使其余 $35$ 个数之和等于这两个数之积。求所选两数之差的绝对值。

\end{problem}

\begin{problem}{练习 25}

两个互不相等的正整数 $a,b$ 满足 $a^2-b^2=2018-2a$，求 $a+b$。

\end{problem}

\Needspace{18\baselineskip}

\subsection{数的性质：整除、质数与完全平方}

\topicline{先分解，再判断整除性；遇到质数或完全平方数条件，还要利用因数的范围。}

\Needspace{12\baselineskip}

\begin{problem}{例题 22 （用整除性筛选）}

将某个正整数 $n$ 代入 $n^5-n$，四个结果中仅一个正确。它是哪一个？A．$7770$；B．$7775$；C．$7776$；D．$7779$。

\end{problem}

\Needspace{6\baselineskip}

\begin{solution}{例题 22 · 讲解}

\think 不必直接算五次方。先分解 $n^5-n$，看它一定能被哪些数整除，再据此排除选项。选出一个结果后，还需找到相应的 $n$。

\solve \[n^5-n=n(n^4-1)=n(n-1)(n+1)(n^2+1).\]
连续三个整数中有一个偶数、一个 $\boxed{3\text{ 的倍数}}$，因此该式必被 $6$ 整除。
先排除奇数选项 B、D。A、C 都是 $6$ 的倍数，还不能仅凭这一点确定答案。
再看因数 $5$：$n$ 除以 $5$ 的余数为 $0,1,2,3,4$ 时，分别有 $n,n-1,n^2+1,n^2+1,n+1$ 被 $5$ 整除。
故该式必被 $30$ 整除，四个选项中只有 A 符合。
取 $n=6$，$6^5-6=7770$，所以 A 确实可以取到。

\insight{整除性可以排除选项；要确认某个值可取到，还应给出对应的整数。}

\end{solution}

\Needspace{12\baselineskip}

\begin{problem}{例题 23 （完全平方数条件）}

自然数 $n$ 满足 $n^2+15n+26$ 是完全平方数，求 $n$。

\end{problem}

\Needspace{6\baselineskip}

\begin{solution}{例题 23 · 讲解}

\think 设原式为 $t^2$。因为一次项系数 $15$ 为奇数，先乘 $4$ 再配方，能避开分数，得到两个整数平方之差为 $121$，再列因数对。

\solve 设 $t$ 是非负整数，且 $t^2=n^2+15n+26$。则
\[(2n+15)^2-(2t)^2=121,\]
即 $(2n+15-2t)(2n+15+2t)=121$。
两因数的积为正，和为 $4n+30>0$，所以均为正整数。又前者不大于后者，因此只有 $(1,121)$、$(11,11)$ 两对。
第一对给出 $2n+15=61$，得 $n=23$、$t=30$；第二对给出 $n=-2$，不符合自然数条件。
故 $n=23$，此时原式为 $900=30^2$。

\insight{一次项系数为奇数时，可先乘 $4$，把配方写成整数形式。}

\end{solution}

\Needspace{12\baselineskip}

\begin{problem}{例题 24 （最大公因数）}

对所有大于 $2$ 的正整数 $n$，求这些整数 $n^5-5n^3+4n$ 的最大公因数。

\end{problem}

\Needspace{6\baselineskip}

\begin{solution}{例题 24 · 讲解}

\think 先提出 $n$，余下部分按 $n^2$ 分解，会得到五个连续整数的积。先证明每个值都是 $120$ 的倍数，再找一个恰好等于 $120$ 的值，才能确定最大公因数。

\solve \[n^5-5n^3+4n=n(n^2-1)(n^2-4)=(n-2)(n-1)n(n+1)(n+2).\]
五个连续整数中至少有一个是 $5$ 的倍数、一个是 $3$ 的倍数。
其乘积也被 $8$ 整除：若 $n$ 为奇数，$n-1,n+1$ 是相邻偶数，其中一个被 $4$ 整除；若 $n$ 为偶数，$n-2,n,n+2$ 中有一个被 $4$ 整除，另有一个偶数提供因数 $2$。
因为 $3,5,8$ 两两互质，乘积被 $120$ 整除。
取 $n=3$，原式恰为 $1\times2\times3\times4\times5=120$，故最大公因数为 $120$。

\insight{求最大公因数，要证明它是公因数，还要说明公因数不能更大。}

\end{solution}

\Needspace{12\baselineskip}

\par\noindent\begin{minipage}{\linewidth}

\begin{problem}{例题 25 （四个连续整数的积加一）}

设 $n$ 为正整数，证明 $(n+2020)(n+2021)(n+2022)(n+2023)+1$ 是完全平方数。

\end{problem}

\Needspace{6\baselineskip}

\begin{solution}[breakable=false]{例题 25 · 讲解}

\think 令 $N=n+2020$，四个因数就成为连续整数。先把首尾一对、中间一对相乘，得到 $N^2+3N$ 与 $N^2+3N+2$。这与例题 5 的两个因子相差 $2$ 的情形相同。

\solve 令 $N=n+2020$，则

\begin{align*}
\text{原式}&=N(N+1)(N+2)(N+3)+1\\
&=[N(N+3)][(N+1)(N+2)]+1\\
&=(N^2+3N)(N^2+3N+2)+1\\
&=(N^2+3N+1)^2.
\end{align*}

因为 $N^2+3N+1$ 是整数，所以原式是完全平方数。更一般地，任意整数 $m$ 都满足

\[
m(m+1)(m+2)(m+3)+1=(m^2+3m+1)^2.
\]

\insight{四个连续整数对应的式子可以两两配对，使两组乘积含有相同的部分；两组乘积相差 $2$，加一后可配成平方。}

\end{solution}

\end{minipage}\par\medskip

\methodcard{整除题可先分解，再找所需的整数因数。质数条件常限制某个正因数为 $1$；完全平方条件可设平方后配成平方差，也可直接把原式配成平方。求最大公因数还应找一个取值来限制它不能更大。}{分别能被 $a,b$ 整除，只有在 $a,b$ 互质时，才能直接推出能被 $ab$ 整除。}

\begin{problem}{练习 26}

已知 $n$ 为正整数，且 $n^4-16n^2+100$ 是质数，求 $n$。

\end{problem}

\begin{problem}{练习 27}

证明 $2018^2+2018^2\times2019^2+2019^2$ 是完全平方数。

\end{problem}

\Needspace{30\baselineskip}

\subsection{几何中的数量关系}

\topicline{把图中的长度与面积写成代数式，再用已学的恒等变形处理。}

\Needspace{12\baselineskip}

\begin{problem}{例题 26 （图形中的面积关系）}

如图，分别以正数 $a,b,m,n$ 为边长作正方形，$m>n$，且 $am-bn=2$、$an+bm=4$。
图①是边长为 $a,b$ 的两个正方形；图②中 $AB=m$、$DC=n$、$BP=m$、$PC=n$，$AB,DC$ 均垂直于 $BC$，阴影为 $\triangle APD$。
\begin{enumerate}[label=(\arabic*)]
\item 若 $a=3,b=4$，求图①涂色面积。
\item 若图①涂色面积为 $3$，四边形 $ABCD$ 面积为 $5$，求图②阴影面积。
\end{enumerate}
\factorfigure

\end{problem}

\Needspace{6\baselineskip}

\begin{solution}{例题 26 · 讲解}

\think 两式平方相加时，含 $abmn$ 的项会抵消，剩下 $(a^2+b^2)(m^2+n^2)$。图①给出 $a^2+b^2$，梯形面积给出 $(m+n)^2$；阴影面积可由梯形减去两个直角三角形求得。

\solve 第 (1) 问：涂色面积为 $a^2+b^2=3^2+4^2=25$。
第 (2) 问：展开并分组得
\begin{align*}
(am-bn)^2+(an+bm)^2&=a^2(m^2+n^2)+b^2(m^2+n^2)\\
&=(a^2+b^2)(m^2+n^2)=20.
\end{align*}
由 $a^2+b^2=3$ 得 $m^2+n^2=\frac{20}{3}$。
梯形两条平行边是 $AB,DC$，它们之间的距离为 $BC=m+n$，故
\[S_{ABCD}=\frac{(m+n)^2}{2}=5.\]
于是 $2mn=(m+n)^2-(m^2+n^2)=10-\frac{20}{3}=\frac{10}{3}$。
又 $S_{ABP}=\frac12m^2$、$S_{PCD}=\frac12n^2$，因此
\[S_{APD}=S_{ABCD}-S_{ABP}-S_{PCD}=5-\frac12\cdot\frac{20}{3}=\frac53.\]

\insight{先写出图中各部分的面积，再看已知的代数关系能求出哪些量。}

\end{solution}

\methodcard{面积关系可以给出平方和、和的平方或乘积。若三边关系中同时有平方项和两两乘积，可尝试配成几个差的平方。}{先确认面积式对应的是图中哪一部分，尤其要找准梯形的高。}

\Needspace{12\baselineskip}

\subsection{选学挑战}

\topicline{这两题需要把前面的方法用于新的情境，可以先试算、记录，再寻找一般的理由。}

\begin{problem}{挑战 1 \starred （黑板上的操作，选学）}

黑板上有 $1,\frac12,\ldots,\frac1{100}$ 共 $100$ 个数。每次删去两个数 $a,b$，写入 $a+b+ab$。经过 $99$ 次操作后，剩下的数是多少？

\end{problem}

\methodcard{操作题可从一次操作前后的等式找不变量。同值条件可先作差，再看所得关系能否用于目标式。}{试算有助于猜想；要说明每次操作或所有取值都成立，还需给出一般理由。}

\begin{problem}{挑战 2 \starred （四次式的同值条件，选学）}

实数 $a\ne b$、正实数 $c$ 满足 $a^4-2019a=b^4-2019b=c$。证明 $-\sqrt c<ab<0$。

\end{problem}

\Needspace{12\baselineskip}
\block{做完本章，回头想一想}
\begin{enumerate}[label=(\arabic*)]
\item 比较例题 16 与练习 17：为什么前者的负值要保留，后者的负值要舍去？依据来自哪里？
\item 比较例题 22 与例题 24：两个五次式分别怎样分解？连续整数的个数相同吗？
\item 从练习 23—25 中选一道，说明每个条件在筛选因数对时用在了哪一步。再合上解析，独立重做一道曾经需要提示的题。
\end{enumerate}

换序自测可按练习 $18,26,21,24,27$ 的顺序进行。每题先判断：题目中的哪条信息提示你采用这种方法？看过解析才做出的题，隔一天再试。


\Needspace{20\baselineskip}
\section*{练习与挑战解析}

\addcontentsline{toc}{section}{练习与挑战解析}

\Needspace{12\baselineskip}

\subsection*{第1章 因式分解的方法拓展}

\Needspace{8\baselineskip}

\block{分组与拆项：让共同因式出现}

\Needspace{7\baselineskip}

\begin{solution}[before skip=4pt plus 1pt,after skip=4pt plus 1pt,top=6pt,bottom=6pt]{练习 1 · 解析}

\think 前两项提出 $a^2$ 后出现 $a+1$；后两项提出负号也出现 $a+1$。提完共同括号后，剩下的式子还可以用平方差公式。

\solve \[a^3+a^2-a-1=a^2(a+1)-(a+1)=(a+1)(a^2-1)=(a+1)^2(a-1).\]

\insight{提取负号要逐项变号；每得到一个因式，都再检查能否继续分解。}

\end{solution}

\Needspace{8\baselineskip}

\block{配方与添项：把结构改成平方差}

\Needspace{7\baselineskip}

\begin{solution}[before skip=4pt plus 1pt,after skip=4pt plus 1pt,top=6pt,bottom=6pt]{练习 2 · 解析}

\think 第一组配成 $(2x-1)^2$ 需要常数 $1$；第二组配成 $(y-2)^2$ 需要常数 $4$。两组相减所产生的常数正好为 $1-4=-3$。

\solve \begin{align*}
4x^2-4x-y^2+4y-3&=(4x^2-4x+1)-(y^2-4y+4)\\
&=(2x-1)^2-(y-2)^2\\
&=(2x-y+1)(2x+y-3).
\end{align*}

\insight{给两组配方时，常数不能各自随意添加；两组总和仍须等于原式。}

\end{solution}

\Needspace{7\baselineskip}

\begin{solution}[before skip=4pt plus 1pt,after skip=4pt plus 1pt,top=6pt,bottom=6pt]{练习 3 · 解析}

\think $a^4+4$ 不能直接使用平方差。补上 $4a^2$ 后可成为 $(a^2+2)^2$，再减去 $(2a)^2$，就出现平方差。

\solve \begin{align*}
a^4+4&=a^4+4a^2+4-4a^2\\
&=(a^2+2)^2-(2a)^2\\
&=(a^2-2a+2)(a^2+2a+2).
\end{align*}
两个二次因式分别为 $(a-1)^2+1$ 与 $(a+1)^2+1$，没有实数根，因而也不能分解为有理系数的一次因式。

\insight{两项平方和不能直接套平方差，但添减同一项后可能形成平方差。}

\end{solution}

\Needspace{8\baselineskip}

\block{整体与换元：减少结构的复杂程度}

\Needspace{7\baselineskip}

\begin{solution}[before skip=4pt plus 1pt,after skip=4pt plus 1pt,top=6pt,bottom=6pt]{练习 4 · 解析}

\think 把常数乘积同为 $6$ 的两对放在一起，可得到 $x^2+7x+6$ 和 $x^2+5x+6$。它们围绕 $x^2+6x+6$ 分别多、少一个 $x$，恰好形成平方差。

\solve \begin{align*}
\text{原式}&=(x^2+7x+6)(x^2+5x+6)+x^2\\
&=[(x^2+6x+6)+x][(x^2+6x+6)-x]+x^2\\
&=(x^2+6x+6)^2-x^2+x^2\\
&=(x^2+6x+6)^2.
\end{align*}

\insight{四个因式不急着展开；优先寻找能产生相同部分的配对。}

\end{solution}

\Needspace{7\baselineskip}

\begin{solution}[before skip=4pt plus 1pt,after skip=4pt plus 1pt,top=6pt,bottom=6pt]{练习 5 · 解析}

\think $x+y$ 与 $xy$ 反复出现。换元后是 $u^2+2uv+v^2-1$，前三项配成平方。代回后的因式里还藏着一个可分解的四项式。

\solve 令 $u=x+y$，$v=xy$，则
\[u(u+2v)+(v+1)(v-1)=(u+v)^2-1=(u+v+1)(u+v-1).\]
因此
\begin{align*}
\text{原式}&=(xy+x+y+1)(xy+x+y-1)\\
&=(x+1)(y+1)(xy+x+y-1).
\end{align*}

\insight{换元后分解完，并不代表原字母下也已分解完；代回后再检查一次。}

\end{solution}

\Needspace{7\baselineskip}

\begin{solution}[before skip=4pt plus 1pt,after skip=4pt plus 1pt,top=6pt,bottom=6pt]{练习 6 · 解析}

\think 换元后将出现 $u^2-2uv+v^2-2u+2v+1$，把 $u-v$ 再看成整体即可配方。最后的平方底数仍然可以分解。

\solve 令 $u=x+y$，$v=xy$，则
\begin{align*}
(u-2v)(u-2)+(v-1)^2&=(u-v)^2-2(u-v)+1\\
&=(u-v-1)^2.
\end{align*}
所以
\begin{align*}
\text{原式}&=(x+y-xy-1)^2\\
&=[-(x-1)(y-1)]^2=(x-1)^2(y-1)^2.
\end{align*}

\insight{完全平方也是积，但平方的底数若还能分解，就必须继续。}

\end{solution}

\Needspace{8\baselineskip}

\block{主元整理：选一个字母建立秩序}

\Needspace{7\baselineskip}

\begin{solution}[before skip=4pt plus 1pt,after skip=4pt plus 1pt,top=6pt,bottom=6pt]{练习 7 · 解析}

\think 关于 $x$ 是三次式，关于 $a$ 是二次式。按 $a$ 整理后，三个系数都含 $x+1$。次数较低只是一个选择线索，真正的依据是整理后能否出现共同因式。

\solve \begin{align*}
\text{原式}&=(x+1)a^2+(2x^2+2x)a+(x^3+x^2-x-1)\\
&=(x+1)a^2+2x(x+1)a+(x+1)(x^2-1)\\
&=(x+1)[a^2+2ax+x^2-1]\\
&=(x+1)[(a+x)^2-1]\\
&=(x+1)(x+a-1)(x+a+1).
\end{align*}

\insight{选择主元后逐项收集系数；不要把“选低次字母”当成必然成功的规则。}

\end{solution}

\Needspace{8\baselineskip}

\block{比较系数：由目标因式倒推}

\Needspace{7\baselineskip}

\begin{solution}[before skip=4pt plus 1pt,after skip=4pt plus 1pt,top=6pt,bottom=6pt]{练习 8 · 解析}

\think 代值为零提供了寻找 $x+1$ 的线索，但不必使用多项式长除法。把 $-2x^2$ 拆成 $x^2-3x^2$，把 $-13x$ 拆成 $-3x-10x$，就能直接分组。

\solve \begin{align*}
x^3-2x^2-13x-10&=(x^3+x^2)+(-3x^2-3x)+(-10x-10)\\
&=(x+1)(x^2-3x-10)\\
&=(x+1)(x+2)(x-5).
\end{align*}

\insight{代值寻找线索，拆项写出依据；不需要把试根和长除法另设成大专题。}

\end{solution}

\Needspace{7\baselineskip}

\begin{solution}[before skip=4pt plus 1pt,after skip=4pt plus 1pt,top=6pt,bottom=6pt]{练习 9 · 解析}

\think 四次式除去二次因式后，剩下二次因式；其首项系数必须为 $2$。设成 $2x^2+px+q$ 后，各次项的系数提供四个等式。

\solve 设
\[2x^4+x^3-ax^2+bx+a+b-1=(x^2+x-6)(2x^2+px+q).\]
比较系数得
\[p+2=1,\quad q+p-12=-a,\quad q-6p=b,\quad -6q=a+b-1.\]
由 $p=-1$ 得 $a=13-q$、$b=q+6$，故 $a+b-1=18$。
于是 $-6q=18$，得 $q=-3$，所以 $a=16$、$b=3$。

\insight{含参数的因式条件，可以转化为有限个系数等式。}

\end{solution}

\Needspace{8\baselineskip}

\block{巩固与选学}

\Needspace{7\baselineskip}

\begin{solution}[before skip=4pt plus 1pt,after skip=4pt plus 1pt,top=6pt,bottom=6pt]{练习 10 · 解析}

\think 前面的三项可以写成 $(x^2+x)^2$，剩余项为 $2(x^2+x)+1$，再把 $x^2+x$ 看成整体。

\solve \begin{align*}
x^4+2x^3+3x^2+2x+1&=(x^2+x)^2+2(x^2+x)+1\\
&=(x^2+x+1)^2.
\end{align*}
$x^2+x+1=(x+\frac12)^2+\frac34>0$，不能再分解为有理系数的一次因式。

\insight{按次数寻找完全平方，再用整体思想检查剩余部分。}

\end{solution}

\Needspace{7\baselineskip}

\begin{solution}[before skip=4pt plus 1pt,after skip=4pt plus 1pt,top=6pt,bottom=6pt]{练习 11 · 解析}

\think 前两项能提出 $x^2y(y-1)$。其余部分要围绕 $y-1$ 改写，其中 $x-xy^2$ 可以用平方差。

\solve \begin{align*}
\text{原式}&=x^2y(y-1)-x(y^2-1)+(y-1)\\
&=(y-1)[x^2y-x(y+1)+1]\\
&=(y-1)[xy(x-1)-(x-1)]\\
&=(x-1)(y-1)(xy-1).
\end{align*}

\insight{第一次分组后若仍复杂，继续围绕已经出现的共同括号调整。}

\end{solution}

\Needspace{7\baselineskip}

\begin{solution}[before skip=4pt plus 1pt,after skip=4pt plus 1pt,top=6pt,bottom=6pt]{挑战 1 · 解析}

\think 记 $A=2x-3y$、$B=3x-2y$，则 $A+B=5(x-y)$。因此原式可写成 $A^3+B^3-(A+B)^3$，展开最后一个立方后会大量抵消。

\solve 由 $(A+B)^3=A^3+3A^2B+3AB^2+B^3$，得
\[A^3+B^3-(A+B)^3=-3AB(A+B).\]
所以原式为
\[-15(x-y)(2x-3y)(3x-2y).\]
这也解释了整体负号和系数 $15$ 的来源。

\insight{出现几个立方时，先比较底数的和差关系，再决定是否展开。}

\end{solution}

\Needspace{12\baselineskip}

\subsection*{第2章 因式分解的应用}

\Needspace{8\baselineskip}

\block{条件求值：把已知整体用起来}

\Needspace{7\baselineskip}

\begin{solution}[before skip=4pt plus 1pt,after skip=4pt plus 1pt,top=6pt,bottom=6pt]{练习 12 · 解析}

\think 三项都含 $ab$；提出后剩下 $a^2+2ab+b^2$，正好由已知的 $a+b$ 决定。

\solve \[a^3b+2a^2b^2+ab^3=ab(a^2+2ab+b^2)=ab(a+b)^2=1\times3^2=9.\]

\insight{先提公因式，能把分散的条件聚到同一个式子里。}

\end{solution}

\Needspace{7\baselineskip}

\begin{solution}[before skip=4pt plus 1pt,after skip=4pt plus 1pt,top=6pt,bottom=6pt]{练习 13 · 解析}

\think 左边加上 $2a^2b^2$ 后成为 $(a^2+b^2)^2$，右边也只含 $a^2+b^2$。这时可把目标整体设为一个新字母。

\solve 移项得 $(a^2+b^2)^2=a^2+b^2+6$。令 $s=a^2+b^2\ge0$，则
\[s^2-s-6=(s-3)(s+2)=0.\]
所以 $s=3$ 或 $s=-2$。由 $s\ge0$ 排除 $-2$，故 $a^2+b^2=3$。

\insight{换元后保留整体自带的范围条件，才能排除不符合原题的解。}

\end{solution}

\Needspace{7\baselineskip}

\begin{solution}[before skip=4pt plus 1pt,after skip=4pt plus 1pt,top=6pt,bottom=6pt]{练习 14 · 解析}

\think 取值相等可以转为差为零。分解后，一个因式是 $m-n$，题目已经说明它不为零，因此另一个因式必须为零。

\solve 两式作差得 $(m-n)(m+n-2)=0$。由 $m\ne n$ 得 $m+n=2$。
于是
\[4=(m+n)^2=m^2+n^2+2mn=5+2mn,\]
所以 $mn=-\frac12$。又
\[a=m(m-2)=m(-n)=-mn=\frac12.\]

\insight{取值相等先作差；非零条件往往用来排除其中一个零因式。}

\end{solution}

\Needspace{8\baselineskip}

\block{简便计算：先分解，再约分}

\Needspace{7\baselineskip}

\begin{solution}[before skip=4pt plus 1pt,after skip=4pt plus 1pt,top=6pt,bottom=6pt]{练习 15 · 解析}

\think 方法拓展章中的添项思想仍可使用。$t^4+64=(t^2+8)^2-(4t)^2$。用 $F(t)$ 表示代数式 $t^2+4t+8$，另一个因式正好是 $F(t-4)$，数字间相差 $4$ 就成为约分依据。

\solve 设 $F(t)=t^2+4t+8$，则
\[t^4+64=(t^2+4t+8)(t^2-4t+8)=F(t)F(t-4).\]
分子各项分解后为
\[F(7)F(3)\,F(15)F(11)\,F(23)F(19)\,F(31)F(27)\,F(39)F(35),\]
分母各项分解后为
\[F(3)F(-1)\,F(11)F(7)\,F(19)F(15)\,F(27)F(23)\,F(35)F(31).\]
因为 $F(t)=(t+2)^2+4>0$，可以约去相同因式，得到
\[\text{原式}=\frac{F(39)}{F(-1)}=\frac{1685}{5}=337.\]

\insight{把两个二次因式写成同一个式子在不同数处的取值，可看清错位约分。}

\end{solution}

\Needspace{8\baselineskip}

\block{平方与非负性：把条件变成限制}

\Needspace{7\baselineskip}

\begin{solution}[before skip=4pt plus 1pt,after skip=4pt plus 1pt,top=6pt,bottom=6pt]{练习 16 · 解析}

\think 为了产生 $12xy$，先用 $(x+6y)^2$，它用掉 $36y^2$；剩下 $16y^2-8y+1$ 又是一个完全平方。

\solve \[x^2+12xy+52y^2-8y+1=(x+6y)^2+(4y-1)^2=0.\]
两个平方均非负，因此 $x+6y=0$、$4y-1=0$，得 $y=\frac14$、$x=-\frac32$。
所以 $x^2-y^2=\frac94-\frac1{16}=\frac{35}{16}$。

\insight{先根据交叉项确定第一个平方，再检查余下的项是否能配成第二个平方。}

\end{solution}

\Needspace{7\baselineskip}

\begin{solution}[before skip=4pt plus 1pt,after skip=4pt plus 1pt,top=6pt,bottom=6pt]{练习 17 · 解析}

\think 两次作差得到同一个二次式 $m^2-mn+n^2$。把它写成平方和，就能判断严格正负，不需要判别式。

\solve \[a-c=m^2-mn+n^2,\qquad c-b=m^2-mn+n^2.\]
而
\[m^2-mn+n^2=\frac12[(m-n)^2+m^2+n^2]>0,\]
其中严格大于零由 $m\ne n$ 保证。因此 $a>c$、$c>b$，即 $b<c<a$。

\insight{比较大小先作差；严格不等式还要说明平方和为何不可能为零。}

\end{solution}

\Needspace{7\baselineskip}

\begin{solution}[before skip=4pt plus 1pt,after skip=4pt plus 1pt,top=6pt,bottom=6pt]{练习 18 · 解析}

\think 配方后两个平方的和不超过 $2$，每个整数的绝对值至多为 $1$。两坐标的取值范围因此很小，可以直接完整列举。

\solve 原不等式等价于 $(x-1)^2+(y-1)^2\le2$。
因为 $x,y$ 是整数，$x-1$、$y-1$ 各只能为 $-1,0,1$，即 $x,y\in\{0,1,2\}$。
这 $3\times3=9$ 种组合的平方和都不超过 $2$，故共有 $9$ 个整数点。

\insight{先用非负性缩小取值范围，再列举全部组合；不必借助圆的知识。}

\end{solution}

\Needspace{8\baselineskip}

\block{整数问题：把等式转成有限因数对}

\Needspace{7\baselineskip}

\begin{solution}[before skip=4pt plus 1pt,after skip=4pt plus 1pt,top=6pt,bottom=6pt]{练习 19 · 解析}

\think 展开该乘积得到 $6ab-9a+10b-15$，利用已知条件可得到积为 $288$。新因式必须分别满足除以 $3$ 的余数和奇偶限制。

\solve 原式等价于 $(3a+5)(2b-3)=288$。记 $p=3a+5$、$q=2b-3$，则 $p-5$ 必须是 $3$ 的倍数，$q$ 为奇数。
因为 $288=2^5\times3^2$，$q$ 只能取 $\pm1,\pm3,\pm9$。
对应的 $p=288/q$ 为 $288,96,32,-288,-96,-32$。
其中只有 $p=32$ 使 $p-5$ 是 $3$ 的倍数，故 $p=32$、$q=9$。
于是 $a=9$、$b=6$，$a+b=15$。

\insight{整数因数对包含负因数；系数带来的余数限制可以有效筛选。}

\end{solution}

\Needspace{7\baselineskip}

\begin{solution}[before skip=4pt plus 1pt,after skip=4pt plus 1pt,top=6pt,bottom=6pt]{练习 20 · 解析}

\think 总和为 $703$。剩余数之和是 $703-a-b$，等于 $ab$ 后可通过加 $1$ 写成 $(a+1)(b+1)$，再用数的范围筛选。

\solve 设所选两数为 $a<b$，则 $1\le a<b\le37$，且
\[703-a-b=ab,\qquad (a+1)(b+1)=704.\]
令 $u=a+1$、$v=b+1$，则 $2\le u<v\le38$。
因 $uv=704$，有 $u\ge\frac{704}{38}>18$，且 $u<\sqrt{704}<27$。
$704=2^6\times11$，在整数 $19$ 至 $26$ 中，只有 $22$ 是其因数。
故 $u=22$、$v=32$，即 $a=21$、$b=31$，两数之差为 $10$。

\insight{列因数对之前先使用范围限制，能减少枚举并保证不漏解。}

\end{solution}

\Needspace{7\baselineskip}

\begin{solution}[before skip=4pt plus 1pt,after skip=4pt plus 1pt,top=6pt,bottom=6pt]{练习 21 · 解析}

\think 移项配方可得到 $(a+1)^2-b^2=2019$。因为 $a+1+b>0$，乘积为正迫使另一因数也为正，只需列正因数对。

\solve 原等式等价于 $(a+1-b)(a+1+b)=2019=3\times673$。
两个因数均为正整数，且前者小于后者，只有 $(1,2019)$ 和 $(3,673)$ 两对。
第一对给出 $a=b=1009$，不符合 $a\ne b$；第二对给出 $a=337$、$b=335$。
所以 $a+b=672$。

\insight{配方把平方差转为整数乘积，再逐一检查正数与不相等条件。}

\end{solution}

\Needspace{8\baselineskip}

\block{数的性质：整除、质数与完全平方}

\Needspace{7\baselineskip}

\begin{solution}[before skip=4pt plus 1pt,after skip=4pt plus 1pt,top=6pt,bottom=6pt]{练习 22 · 解析}

\think 添减 $36n^2$ 后可以用平方差。得到两个正整数因数后，质数条件要求其中一个为 $1$，还要检查最后所得数确为质数。

\solve \begin{align*}
n^4-16n^2+100&=(n^2+10)^2-(6n)^2\\
&=(n^2+6n+10)(n^2-6n+10).
\end{align*}
第一因数 $n^2+6n+10\ge17>1$；第二因数为 $(n-3)^2+1\ge1$。
因此质数条件迫使 $(n-3)^2+1=1$，得 $n=3$。
代回值为 $37$。$37$ 不被 $2,3,5$ 整除，且 $\sqrt{37}<7$，所以它是质数。

\insight{“质数使一个因数为一”必须先说明两个因数都是正整数。}

\end{solution}

\Needspace{7\baselineskip}

\begin{solution}[before skip=4pt plus 1pt,after skip=4pt plus 1pt,top=6pt,bottom=6pt]{练习 23 · 解析}

\think 两数相邻。按 $t$ 展开后，系数依次为 $1,2,3,2,1$，与方法拓展章的完全平方结构一致。

\solve 令 $t=2018$，则
\begin{align*}
t^2+t^2(t+1)^2+(t+1)^2&=t^4+2t^3+3t^2+2t+1\\
&=(t^2+t+1)^2=4074343^2.
\end{align*}
所以原数是完全平方数。又 $t(t+1)$ 为偶数，$t^2+t+1$ 为奇数，所以原数也是奇数。

\insight{相邻数先写成 $t,t+1$，数的规律就变成多项式的结构。}

\end{solution}

\Needspace{8\baselineskip}

\block{选学挑战}

\Needspace{7\baselineskip}

\begin{solution}[before skip=4pt plus 1pt,after skip=4pt plus 1pt,top=6pt,bottom=6pt]{挑战 2 · 解析}

\think 新数加 $1$ 恰好等于两个旧数各加 $1$ 后的乘积。因而所有在场的数各加 $1$ 后再相乘，这个总乘积在每次操作中都保持不变。

\solve 设一次操作写入 $c=a+b+ab$，则 $c+1=(a+1)(b+1)$。
初始时所有数各加 $1$ 后的乘积为
\[2\times\frac32\times\frac43\times\cdots\times\frac{101}{100}=101.\]
每次操作保持这个乘积不变。设最终剩下 $r$，则 $r+1=101$，故 $r=100$。

\insight{操作式含 $a+b+ab$ 时，先看“各加一之后的乘积”是否保持不变。}

\end{solution}

\Needspace{7\baselineskip}

\begin{solution}[before skip=4pt plus 1pt,after skip=4pt plus 1pt,top=6pt,bottom=6pt]{挑战 3 · 解析}

\think 作差后可以消去非零因式 $a-b$，将 $2019$ 表示成关于 $a,b$ 的积。代入 $c$ 后，四次项抵消，留下能直接判断符号的因式。

\solve 由 $a^4-b^4=2019(a-b)$，且 $a\ne b$，得
\[(a+b)(a^2+b^2)=2019.\]
所以 $a+b>0$，并且
\begin{align*}
c&=a^4-a(a+b)(a^2+b^2)\\
&=-ab(a^2+ab+b^2).
\end{align*}
因为 $a^2+ab+b^2=\frac12[(a+b)^2+a^2+b^2]>0$，而 $c>0$，故 $ab<0$。
再由
\[c-(ab)^2=-ab(a+b)^2>0,\]
得 $(ab)^2<c$。结合 $ab<0$，得到 $-\sqrt c<ab<0$。

\insight{同值条件先作差；把得到的低次关系代回原式，能把高次问题转成符号判断。}

\end{solution}

\end{document}
